\subsection{Свойства функций, имеющий пределы: единственность предела, предел модуля, предельные переходы в неравенствах и теорема о зажатой функции (о трёх функциях)}

\textbf{Свойства функций, имеющих предел}\\
Всюду далее считаем, что пределы рассматриваются при $x \to x_0$ (или $x \to \infty$).

\begin{tcolorbox}[title=1. Единственность предела]
	\textbf{Теорема.} Если функция $f(x)$ имеет предел, то он единственен.
	
	\textbf{Доказательство (от противного):}
	\begin{enumerate}
		\item Допустим, $\lim f(x) = A$ и $\lim f(x) = B$, причем $A \neq B$.
		\item Выберем $\varepsilon = \frac{|A - B|}{2} > 0$ (половина расстояния между точками).
		\item По определению предела:
		\begin{itemize}
			\item $\exists \delta_1: \forall x \in \dot{U}_{\delta_1}(x_0) \Rightarrow f(x) \in (A-\varepsilon, A+\varepsilon)$
			\item $\exists \delta_2: \forall x \in \dot{U}_{\delta_2}(x_0) \Rightarrow f(x) \in (B-\varepsilon, B+\varepsilon)$
		\end{itemize}
		\item Возьмем $\delta = \min(\delta_1, \delta_2)$. Тогда для любого $x \in \dot{U}_\delta(x_0)$ значение $f(x)$ должно находиться одновременно в $\varepsilon$-окрестности $A$ и $B$.
		\item Но эти окрестности не пересекаются (так как радиус меньше половины расстояния). \textbf{Противоречие.}
	\end{enumerate}
\end{tcolorbox}

\begin{tcolorbox}[title=2. Предел модуля]
	\textbf{Теорема.} Если $\lim f(x) = A$, то $\lim |f(x)| = |A|$.
	
	\textbf{Доказательство:}
	Используем неравенство треугольника вида $| |a| - |b| | \leq |a - b|$.
	\begin{itemize}
		\item Зададим $\forall \varepsilon > 0$.
		\item Так как $f(x) \to A$, то $\exists \delta: \forall x \in \dot{U}_\delta \Rightarrow |f(x) - A| < \varepsilon$.
		\item Тогда:
		$$ ||f(x)| - |A|| \leq |f(x) - A| < \varepsilon $$
		\item Это и означает, что $|f(x)| \to |A|$.
	\end{itemize}
	\textit{Важно: Обратное верно только если $A=0$.}
\end{tcolorbox}

\begin{tcolorbox}[title=3. Предельный переход в неравенствах]
	\textbf{Теорема.} Пусть $\lim f(x) = A$ и $\lim g(x) = B$.
	Если в некоторой окрестности точки $x_0$ выполняется $f(x) \leq g(x)$, то $A \leq B$.
	
	\textbf{Доказательство (от противного):}
	\begin{enumerate}
		\item Допустим, что $A > B$.
		\item Возьмем $\varepsilon = \frac{A - B}{2} > 0$.
		\item Найдется окрестность, где:
		$$ |f(x) - A| < \varepsilon \Rightarrow f(x) > A - \varepsilon = A - \frac{A-B}{2} = \frac{A+B}{2} $$
		$$ |g(x) - B| < \varepsilon \Rightarrow g(x) < B + \varepsilon = B + \frac{A-B}{2} = \frac{A+B}{2} $$
		\item Получаем $g(x) < \frac{A+B}{2} < f(x)$, то есть $g(x) < f(x)$.
		\item Это противоречит условию $f(x) \leq g(x)$. Значит, предположение неверно, и $A \leq B$.
	\end{enumerate}
	\textit{Замечание: Строгое неравенство $f(x) < g(x)$ при переходе к пределу может стать нестрогим (например, $1/x^2 > 0$, но предел равен 0).}
\end{tcolorbox}

\begin{tcolorbox}[title=4. Теорема о зажатой функции (о двух милиционерах)]
	\textbf{Теорема.} Пусть в некоторой окрестности точки $x_0$ выполняется:
	$$ g(x) \leq f(x) \leq h(x) $$
	И пусть $\lim_{x\to x_0} g(x) = \lim_{x\to x_0} h(x) = A$.
	Тогда $\lim_{x\to x_0} f(x) = A$.
	
	\textbf{Доказательство:}
	\begin{enumerate}
		\item Запишем определение предела для крайних функций. $\forall \varepsilon > 0 \exists \delta$:
		\begin{itemize}
			\item $A - \varepsilon < g(x) < A + \varepsilon$
			\item $A - \varepsilon < h(x) < A + \varepsilon$
		\end{itemize}
		\item Используем двойное неравенство из условия:
		$$ A - \varepsilon < g(x) \leq f(x) \leq h(x) < A + \varepsilon $$
		\item Уберем промежуточные члены:
		$$ A - \varepsilon < f(x) < A + \varepsilon $$
		$$ |f(x) - A| < \varepsilon $$
		\item Это и означает по определению, что $\lim f(x) = A$.
	\end{enumerate}
\end{tcolorbox}
